\documentclass[11pt,a4paper]{article}

%====================================================
% PACKAGES
%====================================================
\usepackage[utf8]{inputenc}
\usepackage[T1]{fontenc}
\usepackage[french]{babel}
\usepackage[a4paper,margin=2cm]{geometry}
\usepackage{amsmath,amssymb,amsthm,mathtools}
\usepackage{physics}
\usepackage{siunitx}
\usepackage{lmodern}
\usepackage{xcolor}
\usepackage{hyperref}
\usepackage{fancyhdr}
\usepackage{lastpage}
\usepackage{enumitem}
\usepackage{array}
\usepackage{tabularx}
\usepackage{graphicx}
\usepackage{tikz}
\usepackage{pgfplots}
\usepackage{multicol}
\pgfplotsset{compat=1.18}

\sisetup{locale=FR}

%====================================================
% MISE EN PAGE
%====================================================
\hypersetup{
    colorlinks=true,
    linkcolor=blue!50!black,
    urlcolor=blue!50!black
}

\setlength{\headheight}{14pt}
\pagestyle{fancy}
\fancyhf{}
\fancyhead[L]{Terminale générale -- Spécialité Physique-Chimie}
\fancyhead[C]{Cours complet -- Mécanique}
\fancyhead[R]{Page \thepage/\pageref{LastPage}}
\fancyfoot[C]{Cours détaillé avec démonstrations, méthodes et clés de compréhension}

%====================================================
% ENVIRONNEMENTS
%====================================================
\theoremstyle{definition}
\newtheorem{definition}{Définition}[section]
\newtheorem{exemple}[definition]{Exemple}
\newtheorem{remarque}[definition]{Remarque}
\newtheorem{methode}[definition]{Méthode}
\newtheorem{interpretation}[definition]{Clé de compréhension}
\newtheorem{vocabulaire}[definition]{Vocabulaire}

\theoremstyle{plain}
\newtheorem{propriete}[definition]{Propriété}
\newtheorem{theoreme}[definition]{Théorème}
\newtheorem{proposition}[definition]{Proposition}
\newtheorem{lemme}[definition]{Lemme}
\newtheorem{corollaire}[definition]{Corollaire}

\theoremstyle{remark}
\newtheorem*{demonstrationlibre}{Démonstration}
\newtheorem*{conclusion}{Conclusion}

%====================================================
% COMMANDES
%====================================================
\newcommand{\R}{\mathbb{R}}
\newcommand{\vect}[1]{\overrightarrow{#1}}
\newcommand{\uu}{\vec{u}}
\newcommand{\vvv}{\vec{v}}
\newcommand{\aaa}{\vec{a}}
\newcommand{\OM}{\overrightarrow{OM}}
\newcommand{\e}{\mathrm{e}}

\begin{document}

\begin{center}
    {\LARGE \textbf{Cours de Terminale spécialité Physique-Chimie}}\\[0.2cm]
    {\Huge \textbf{Mécanique -- Mouvement et interactions}}\\[0.3cm]
    {\large Version très complète avec démonstrations, méthodes, exemples et interprétations physiques}
\end{center}

\vspace{0.4cm}

\begin{center}
\renewcommand{\arraystretch}{1.4}
\begin{tabular}{|p{0.9\textwidth}|}
\hline
\textbf{Objectifs du chapitre}\\
Comprendre comment décrire un mouvement, comment relier ce mouvement aux forces exercées sur un système, comment utiliser les lois de Newton, comment modéliser des mouvements usuels (chute libre, mouvement circulaire, projectile) et comment interpréter correctement les grandeurs mécaniques fondamentales : position, vitesse, accélération, masse, force, quantité de mouvement.\\
\hline
\end{tabular}
\end{center}

\tableofcontents
\newpage

%====================================================
\section{Introduction générale : que cherche la mécanique ?}
%====================================================

\begin{definition}
La \textbf{mécanique} est la partie de la physique qui étudie les mouvements des systèmes matériels et les causes de ces mouvements.
\end{definition}

\begin{interpretation}
La mécanique répond à deux grandes questions :
\begin{itemize}
    \item \textbf{Comment} le système se déplace-t-il ? C'est la \textbf{cinématique}.
    \item \textbf{Pourquoi} se déplace-t-il ainsi ? C'est la \textbf{dynamique}.
\end{itemize}
\end{interpretation}

\begin{remarque}
En Terminale, on sépare donc très souvent :
\begin{itemize}
    \item la description du mouvement ;
    \item l'étude des forces ;
    \item l'application des lois de Newton ;
    \item l'étude de cas types.
\end{itemize}
\end{remarque}

%====================================================
\section{Système, référentiel, trajectoire}
%====================================================

\subsection{Système étudié}

\begin{definition}
Le \textbf{système} est l'objet, ou l'ensemble d'objets, que l'on choisit d'étudier.
\end{definition}

\begin{exemple}
Le système peut être :
\begin{itemize}
    \item une balle ;
    \item une voiture ;
    \item un skieur ;
    \item une planète ;
    \item un électron ;
    \item une fusée.
\end{itemize}
\end{exemple}

\begin{remarque}
Avant toute étude mécanique, il faut toujours identifier clairement le système.
\end{remarque}

\subsection{Référentiel}

\begin{definition}
Un \textbf{référentiel} est un solide de référence muni :
\begin{itemize}
    \item d'un repère d'espace ;
    \item d'une horloge.
\end{itemize}
Il permet de décrire la position et le mouvement d'un système.
\end{definition}

\begin{exemple}
On utilise souvent :
\begin{itemize}
    \item le référentiel terrestre ;
    \item le référentiel géocentrique ;
    \item le référentiel héliocentrique.
\end{itemize}
\end{exemple}

\begin{interpretation}
Un mouvement n'existe jamais « en soi » : il est toujours décrit \textbf{par rapport à un référentiel}.  

Un même objet peut être immobile dans un référentiel et en mouvement dans un autre.
\end{interpretation}

\subsection{Trajectoire}

\begin{definition}
La \textbf{trajectoire} est l'ensemble des positions successives occupées par un point du système au cours du temps.
\end{definition}

\begin{definition}
Selon sa forme, on distingue :
\begin{itemize}
    \item une trajectoire \textbf{rectiligne} ;
    \item une trajectoire \textbf{circulaire} ;
    \item une trajectoire \textbf{curviligne}.
\end{itemize}
\end{definition}

\begin{exemple}
\begin{itemize}
    \item Une voiture sur une ligne droite : trajectoire rectiligne.
    \item Une nacelle de grande roue : trajectoire circulaire.
    \item Un ballon lancé : trajectoire curviligne.
\end{itemize}
\end{exemple}

%====================================================
\section{Position, vecteur position, vitesse, accélération}
%====================================================

\subsection{Vecteur position}

\begin{definition}
Dans un repère donné, la position du point matériel $M$ est repérée par le \textbf{vecteur position}
\[
\OM.
\]
\end{definition}

Si le point a pour coordonnées $(x(t),y(t),z(t))$, alors
\[
\OM(t)=x(t)\,\vec{i}+y(t)\,\vec{j}+z(t)\,\vec{k}.
\]

\subsection{Vecteur vitesse moyenne}

\begin{definition}
Entre deux instants $t_1$ et $t_2$, la \textbf{vitesse moyenne} est définie par
\[
\vec{v}_{\text{moy}}=\frac{\overrightarrow{M_1M_2}}{t_2-t_1}.
\]
\end{definition}

\begin{remarque}
La vitesse moyenne dépend d'un intervalle de temps. Elle ne décrit pas précisément le mouvement à un instant donné.
\end{remarque}

\subsection{Vecteur vitesse instantanée}

\begin{definition}
Le \textbf{vecteur vitesse instantanée} du point $M$ à l'instant $t$ est la dérivée temporelle du vecteur position :
\[
\vec{v}(t)=\dv{\OM}{t}.
\]
\end{definition}

Si
\[
\OM(t)=x(t)\,\vec{i}+y(t)\,\vec{j}+z(t)\,\vec{k},
\]
alors
\[
\vec{v}(t)=x'(t)\,\vec{i}+y'(t)\,\vec{j}+z'(t)\,\vec{k}.
\]

\begin{propriete}
Le vecteur vitesse instantanée est tangent à la trajectoire et orienté dans le sens du mouvement.
\end{propriete}

\begin{interpretation}
La direction de la vitesse indique la direction vers laquelle l'objet « part immédiatement ».  

La valeur
\[
v=\norm{\vec{v}}
\]
s'appelle la \textbf{vitesse} et s'exprime en \si{m.s^{-1}}.
\end{interpretation}

\subsection{Vecteur accélération}

\begin{definition}
Le \textbf{vecteur accélération} est la dérivée temporelle du vecteur vitesse :
\[
\vec{a}(t)=\dv{\vec{v}}{t}=\dv[2]{\OM}{t}.
\]
\end{definition}

Si
\[
\vec{v}(t)=v_x(t)\,\vec{i}+v_y(t)\,\vec{j}+v_z(t)\,\vec{k},
\]
alors
\[
\vec{a}(t)=\dv{v_x}{t}\,\vec{i}+\dv{v_y}{t}\,\vec{j}+\dv{v_z}{t}\,\vec{k}.
\]

\begin{interpretation}
L'accélération mesure la manière dont le vecteur vitesse varie :
\begin{itemize}
    \item variation de la \textbf{valeur} de la vitesse ;
    \item variation de la \textbf{direction} de la vitesse ;
    \item ou les deux.
\end{itemize}
\end{interpretation}

\begin{exemple}
Même à vitesse constante, un mouvement circulaire possède une accélération, car la direction de la vitesse change.
\end{exemple}

%====================================================
\section{Mouvements usuels}
%====================================================

\subsection{Mouvement rectiligne uniforme}

\begin{definition}
Un mouvement est \textbf{rectiligne uniforme} si :
\begin{itemize}
    \item la trajectoire est une droite ;
    \item la vitesse est constante.
\end{itemize}
\end{definition}

\begin{propriete}
Dans un mouvement rectiligne uniforme,
\[
\vec{a}=\vec{0}.
\]
\end{propriete}

\begin{demonstrationlibre}
Si la vitesse est constante, alors le vecteur vitesse ne varie pas au cours du temps. Sa dérivée temporelle est donc nulle :
\[
\vec{a}=\dv{\vec{v}}{t}=\vec{0}.
\]
\end{demonstrationlibre}

\begin{exemple}
Une voiture roulant en ligne droite à vitesse strictement constante sur une route horizontale peut être modélisée par un mouvement rectiligne uniforme.
\end{exemple}

\subsection{Mouvement rectiligne uniformément varié}

\begin{definition}
Un mouvement est \textbf{rectiligne uniformément varié} si la trajectoire est rectiligne et si l'accélération est constante.
\end{definition}

\begin{propriete}
Si l'accélération est constante dans un mouvement rectiligne, alors :
\[
v(t)=v_0+at
\]
et
\[
x(t)=x_0+v_0t+\frac12 at^2.
\]
\end{propriete}

\begin{demonstrationlibre}
Dans un mouvement rectiligne, on note l'accélération scalaire $a$ constante.  

Comme
\[
a=\dv{v}{t},
\]
on intègre :
\[
v(t)=v_0+at.
\]

Ensuite,
\[
v(t)=\dv{x}{t}=v_0+at.
\]
On intègre à nouveau :
\[
x(t)=x_0+v_0t+\frac12 at^2.
\]
\end{demonstrationlibre}

\begin{interpretation}
Le terme en $t^2$ apparaît dès qu'il y a accélération constante. C'est la signature mathématique d'un mouvement uniformément varié.
\end{interpretation}

\subsection{Mouvement circulaire uniforme}

\begin{definition}
Un mouvement est \textbf{circulaire uniforme} si la trajectoire est un cercle et si la valeur de la vitesse est constante.
\end{definition}

\begin{propriete}
Dans un mouvement circulaire uniforme :
\begin{itemize}
    \item la norme de la vitesse est constante ;
    \item le vecteur vitesse change de direction ;
    \item l'accélération est dirigée vers le centre du cercle.
\end{itemize}
\end{propriete}

\begin{definition}
On appelle \textbf{accélération centripète} l'accélération d'un mouvement circulaire uniforme :
\[
a=\frac{v^2}{R}
\]
où $R$ est le rayon du cercle.
\end{definition}

\begin{interpretation}
Même si « l'objet ne va pas plus vite », il accélère, car il tourne.  
L'accélération mesure aussi le changement de direction.
\end{interpretation}

%====================================================
\section{Interactions et forces}
%====================================================

\subsection{Notion d'interaction}

\begin{definition}
Une \textbf{interaction} est une action réciproque entre deux systèmes.
\end{definition}

\begin{definition}
Une interaction peut être :
\begin{itemize}
    \item \textbf{de contact} : support, tension, frottement, poussée, réaction ;
    \item \textbf{à distance} : gravitation, interaction électrostatique, interaction magnétique.
\end{itemize}
\end{definition}

\subsection{Force}

\begin{definition}
Une \textbf{force} exercée sur un système est modélisée par un vecteur caractérisé par :
\begin{itemize}
    \item son point d'application ;
    \item sa direction ;
    \item son sens ;
    \item sa valeur.
\end{itemize}
\end{definition}

\begin{definition}
L'unité de la force est le \textbf{newton} (\si{N}).
\end{definition}

\begin{remarque}
Une force n'est pas « une réserve de mouvement ». Une force traduit l'action d'un système sur un autre.
\end{remarque}

\subsection{Quelques forces usuelles}

\begin{proposition}
Les forces fréquemment rencontrées en Terminale sont :
\begin{itemize}
    \item le \textbf{poids} ;
    \item la \textbf{réaction du support} ;
    \item la \textbf{tension} d'un fil ;
    \item la \textbf{force de frottement} ;
    \item la \textbf{force électrique} ;
    \item la \textbf{force gravitationnelle}.
\end{itemize}
\end{proposition}

\subsection{Poids}

\begin{definition}
Le \textbf{poids} d'un système de masse $m$ dans un champ de pesanteur $\vec{g}$ est :
\[
\vec{P}=m\vec{g}.
\]
\end{definition}

\begin{remarque}
Près de la surface terrestre, on prend souvent
\[
g\simeq \SI{9,81}{m.s^{-2}}
\]
ou parfois
\[
g\simeq \SI{9,8}{m.s^{-2}}.
\]
\end{remarque}

\begin{interpretation}
Le poids n'est pas la masse.  
\begin{itemize}
    \item La \textbf{masse} mesure l'inertie du système, en \si{kg}.
    \item Le \textbf{poids} est une force, en \si{N}.
\end{itemize}
\end{interpretation}

\subsection{Réaction du support}

\begin{definition}
La \textbf{réaction du support} est la force exercée par un support sur le système en contact avec lui.
\end{definition}

\begin{remarque}
Elle peut se décomposer en :
\begin{itemize}
    \item une composante normale au support ;
    \item une composante tangentielle liée au frottement.
\end{itemize}
\end{remarque}

\subsection{Tension d'un fil}

\begin{definition}
La \textbf{tension} est la force exercée par un fil, une corde ou un câble tendu sur le système attaché.
\end{definition}

\subsection{Frottements}

\begin{definition}
Les \textbf{forces de frottement} s'opposent au mouvement ou à la tendance au mouvement.
\end{definition}

\begin{remarque}
En Terminale, on modélise souvent les frottements de manière simple :
\begin{itemize}
    \item frottements négligeables ;
    \item force de frottement constante ;
    \item force opposée à la vitesse.
\end{itemize}
\end{remarque}

%====================================================
\section{Principe d'inertie}
%====================================================

\begin{theoreme}[Première loi de Newton]
Dans un référentiel galiléen, si la résultante des forces extérieures exercées sur un système est nulle, alors le système :
\begin{itemize}
    \item reste au repos s'il était au repos ;
    \item conserve un mouvement rectiligne uniforme s'il était en mouvement.
\end{itemize}
\end{theoreme}

\begin{remarque}
Un \textbf{référentiel galiléen} est un référentiel dans lequel le principe d'inertie est vérifié. À l'échelle usuelle, le référentiel terrestre est souvent assimilé à un référentiel galiléen.
\end{remarque}

\begin{interpretation}
Le mouvement uniforme n'a pas besoin d'être « entretenu » par une force résultante.  

Ce qui nécessite une force résultante non nulle, c'est la \textbf{modification} du vecteur vitesse.
\end{interpretation}

\begin{propriete}
Si
\[
\sum \vec{F}_{\text{ext}}=\vec{0},
\]
alors
\[
\vec{a}=\vec{0}.
\]
\end{propriete}

\begin{remarque}
En pratique :
\begin{itemize}
    \item si un objet est immobile ou en mouvement rectiligne uniforme, la somme des forces peut être nulle ;
    \item réciproquement, si la somme des forces est nulle, alors l'accélération est nulle.
\end{itemize}
\end{remarque}

%====================================================
\section{Deuxième loi de Newton}
%====================================================

\subsection{Énoncé}

\begin{theoreme}[Deuxième loi de Newton]
Dans un référentiel galiléen, la somme des forces extérieures exercées sur un système de masse constante $m$ vérifie
\[
\sum \vec{F}_{\text{ext}}=m\vec{a}.
\]
\end{theoreme}

\begin{interpretation}
Cette relation est l'équation fondamentale de la dynamique.

Elle relie :
\begin{itemize}
    \item la \textbf{cause} : les forces ;
    \item la \textbf{conséquence} : l'accélération.
\end{itemize}
\end{interpretation}

\subsection{Sens physique}

\begin{interpretation}
Cette loi montre que :
\begin{itemize}
    \item si la résultante des forces est nulle, l'accélération est nulle ;
    \item plus la force résultante est grande, plus l'accélération est grande ;
    \item plus la masse est grande, plus il est difficile de modifier le mouvement.
\end{itemize}
\end{interpretation}

\begin{remarque}
La masse traduit l'\textbf{inertie} du système : sa résistance à la modification de son mouvement.
\end{remarque}

\subsection{Projection dans un repère}

\begin{propriete}
Dans un repère cartésien,
\[
\sum \vec{F}=m\vec{a}
\]
équivaut à trois équations scalaires :
\[
\sum F_x=ma_x,\qquad \sum F_y=ma_y,\qquad \sum F_z=ma_z.
\]
\end{propriete}

\begin{methode}
Pour exploiter la deuxième loi de Newton :
\begin{enumerate}
    \item choisir le système ;
    \item choisir le référentiel ;
    \item faire le bilan des forces ;
    \item choisir un repère adapté ;
    \item projeter l'équation vectorielle sur les axes.
\end{enumerate}
\end{methode}

%====================================================
\section{Troisième loi de Newton}
%====================================================

\begin{theoreme}[Troisième loi de Newton -- action-réaction]
Si un système $A$ exerce sur un système $B$ une force $\vec{F}_{A\to B}$, alors le système $B$ exerce sur $A$ une force
\[
\vec{F}_{B\to A}=-\vec{F}_{A\to B}.
\]
\end{theoreme}

\begin{interpretation}
Les forces d'interaction apparaissent toujours par paires :
\begin{itemize}
    \item même direction ;
    \item même valeur ;
    \item sens opposés ;
    \item appliquées sur deux systèmes différents.
\end{itemize}
\end{interpretation}

\begin{remarque}
L'erreur classique est de croire que ces deux forces « se compensent » sur un même système. C'est faux : elles s'appliquent à deux objets différents.
\end{remarque}

%====================================================
\section{Quantité de mouvement}
%====================================================

\subsection{Définition}

\begin{definition}
La \textbf{quantité de mouvement} d'un système de masse $m$ animé d'une vitesse $\vec{v}$ est :
\[
\vec{p}=m\vec{v}.
\]
\end{definition}

\begin{definition}
L'unité de la quantité de mouvement est
\[
\si{kg.m.s^{-1}}.
\]
\end{definition}

\subsection{Lien avec la deuxième loi de Newton}

\begin{propriete}
Pour un système de masse constante,
\[
\sum \vec{F}_{\text{ext}}=\dv{\vec{p}}{t}.
\]
\end{propriete}

\begin{demonstrationlibre}
On a
\[
\vec{p}=m\vec{v}.
\]
Si la masse est constante,
\[
\dv{\vec{p}}{t}=m\dv{\vec{v}}{t}=m\vec{a}.
\]
Or, d'après la deuxième loi de Newton,
\[
m\vec{a}=\sum \vec{F}_{\text{ext}}.
\]
Donc
\[
\sum \vec{F}_{\text{ext}}=\dv{\vec{p}}{t}.
\]
\end{demonstrationlibre}

\begin{interpretation}
La force résultante mesure la rapidité de variation de la quantité de mouvement.
\end{interpretation}

%====================================================
\section{Mouvement dans un champ de pesanteur uniforme}
%====================================================

\subsection{Chute libre}

\begin{definition}
On parle de \textbf{chute libre} lorsqu'un système est soumis uniquement à son poids.
\end{definition}

\begin{remarque}
En pratique, on néglige les frottements de l'air.
\end{remarque}

\begin{propriete}
En chute libre, dans un champ de pesanteur uniforme,
\[
\vec{a}=\vec{g}.
\]
\end{propriete}

\begin{demonstrationlibre}
Le système n'est soumis qu'à son poids :
\[
\sum \vec{F}=\vec{P}=m\vec{g}.
\]
D'après la deuxième loi de Newton,
\[
m\vec{a}=m\vec{g}.
\]
Comme $m\neq0$,
\[
\vec{a}=\vec{g}.
\]
\end{demonstrationlibre}

\subsection{Équations horaires de la chute verticale}

On choisit un axe vertical orienté vers le haut. Alors
\[
a_z=-g.
\]

\begin{propriete}
Pour une chute verticale,
\[
v_z(t)=v_{z0}-gt
\]
et
\[
z(t)=z_0+v_{z0}t-\frac12 gt^2.
\]
\end{propriete}

\begin{demonstrationlibre}
On part de
\[
a_z=\dv{v_z}{t}=-g.
\]
En intégrant :
\[
v_z(t)=v_{z0}-gt.
\]
Puis
\[
v_z=\dv{z}{t}.
\]
On intègre :
\[
z(t)=z_0+v_{z0}t-\frac12 gt^2.
\]
\end{demonstrationlibre}

\subsection{Cas du lancer horizontal}

\begin{propriete}
Pour un lancer horizontal sans frottements :
\begin{itemize}
    \item sur l'axe horizontal :
    \[
    a_x=0,\qquad v_x=v_0,\qquad x(t)=x_0+v_0t ;
    \]
    \item sur l'axe vertical :
    \[
    a_z=-g,\qquad v_z=-gt,\qquad z(t)=z_0-\frac12 gt^2.
    \]
\end{itemize}
\end{propriete}

\begin{interpretation}
Le mouvement horizontal et le mouvement vertical sont indépendants :
\begin{itemize}
    \item horizontalement : mouvement uniforme ;
    \item verticalement : chute libre.
\end{itemize}
\end{interpretation}

\subsection{Projectile lancé avec une vitesse initiale quelconque}

Si
\[
\vec{v}_0=v_{0x}\vec{i}+v_{0z}\vec{k},
\]
alors
\[
x(t)=x_0+v_{0x}t,\qquad z(t)=z_0+v_{0z}t-\frac12 gt^2.
\]

\begin{propriete}
La trajectoire d'un projectile sans frottements, dans un champ de pesanteur uniforme, est une parabole.
\end{propriete}

\begin{demonstrationlibre}
On a
\[
x(t)=x_0+v_{0x}t.
\]
Donc
\[
t=\frac{x-x_0}{v_{0x}}
\]
si $v_{0x}\neq 0$.

En remplaçant dans l'expression de $z(t)$ :
\[
z=z_0+v_{0z}\frac{x-x_0}{v_{0x}}-\frac12 g\left(\frac{x-x_0}{v_{0x}}\right)^2.
\]
Cette expression est du second degré en $x$, donc la trajectoire est une parabole.
\end{demonstrationlibre}

%====================================================
\section{Mouvement circulaire et force centripète}
%====================================================

\subsection{Accélération centripète}

\begin{propriete}
Dans un mouvement circulaire uniforme de rayon $R$ et de vitesse de valeur $v$,
\[
a=\frac{v^2}{R}.
\]
\end{propriete}

\begin{remarque}
Cette accélération est dirigée vers le centre du cercle.
\end{remarque}

\subsection{Conséquence dynamique}

\begin{propriete}
Pour maintenir un mouvement circulaire uniforme, la résultante des forces doit être dirigée vers le centre et vérifier
\[
\sum F_{\text{radiale}}=m\frac{v^2}{R}.
\]
\end{propriete}

\begin{interpretation}
Le rôle de la force résultante n'est pas ici d'augmenter la vitesse, mais de courber la trajectoire.
\end{interpretation}

\begin{exemple}
Pour une voiture dans un virage, l'adhérence fournit la force latérale nécessaire au mouvement circulaire.
\end{exemple}

%====================================================
\section{Mouvement dans un champ électrique uniforme}
%====================================================

\begin{definition}
Une charge électrique $q$ placée dans un champ électrique uniforme $\vec{E}$ subit une force :
\[
\vec{F}=q\vec{E}.
\]
\end{definition}

\begin{propriete}
Dans un champ électrique uniforme, si les autres forces sont négligées,
\[
m\vec{a}=q\vec{E}
\qquad \Longrightarrow \qquad
\vec{a}=\frac{q}{m}\vec{E}.
\]
\end{propriete}

\begin{interpretation}
L'accélération dépend :
\begin{itemize}
    \item du champ $\vec{E}$ ;
    \item de la charge $q$ ;
    \item de la masse $m$.
\end{itemize}
\end{interpretation}

\begin{remarque}
Si $q>0$, la force est dans le sens de $\vec{E}$.  
Si $q<0$, elle est de sens opposé.
\end{remarque}

%====================================================
\section{Méthode générale de résolution en mécanique}
%====================================================

\begin{methode}
Pour résoudre correctement un exercice de mécanique, on suit presque toujours ce plan :
\begin{enumerate}
    \item \textbf{Définir le système.}
    \item \textbf{Choisir le référentiel}, en précisant s'il est supposé galiléen.
    \item \textbf{Faire le bilan des forces} extérieures.
    \item \textbf{Choisir un repère} adapté à la situation.
    \item \textbf{Appliquer la loi de Newton} :
    \[
    \sum \vec{F}=m\vec{a}.
    \]
    \item \textbf{Projeter} sur les axes.
    \item \textbf{Résoudre les équations} obtenues.
    \item \textbf{Interpréter physiquement} le résultat.
\end{enumerate}
\end{methode}

\begin{remarque}
Le choix du repère est stratégique :
\begin{itemize}
    \item axe horizontal / vertical pour les chutes et projectiles ;
    \item axe tangent / normal pour les trajectoires courbes ;
    \item axe parallèle au plan incliné pour les glissements sur pente.
\end{itemize}
\end{remarque}

%====================================================
\section{Exemples fondamentaux détaillés}
%====================================================

\subsection{Exemple 1 : solide immobile sur une table}

\begin{exemple}
Un objet de masse $m$ est immobile sur une table horizontale.
\end{exemple}

\textbf{Système :} l'objet.  

\textbf{Forces :}
\begin{itemize}
    \item le poids $\vec{P}=m\vec{g}$, vers le bas ;
    \item la réaction du support $\vec{R}$, vers le haut.
\end{itemize}

Comme l'objet est immobile,
\[
\vec{a}=\vec{0}.
\]
Donc
\[
\vec{R}+\vec{P}=\vec{0}.
\]
Ainsi,
\[
R=P=mg.
\]

\begin{interpretation}
Le support n'annule pas le poids ; il exerce une force opposée qui équilibre le système.
\end{interpretation}

\subsection{Exemple 2 : chute libre verticale}

\begin{exemple}
On lâche une bille sans vitesse initiale depuis une hauteur $h$.
\end{exemple}

Le seul effort est le poids :
\[
m\vec{a}=m\vec{g}.
\]
Sur l'axe vertical orienté vers le haut :
\[
a_z=-g.
\]
Comme la vitesse initiale est nulle :
\[
v_z(t)=-gt,
\qquad
z(t)=h-\frac12 gt^2.
\]

Le sol est atteint lorsque
\[
z(t)=0.
\]
Donc
\[
h-\frac12 gt^2=0
\quad \Longrightarrow \quad
t=\sqrt{\frac{2h}{g}}.
\]

\subsection{Exemple 3 : lancer horizontal}

\begin{exemple}
Une balle est lancée horizontalement depuis une hauteur $h$ avec une vitesse initiale $v_0$.
\end{exemple}

Équations horaires :
\[
x(t)=v_0t,
\qquad
z(t)=h-\frac12 gt^2.
\]

Le temps de chute est donné par
\[
h-\frac12 gt^2=0
\quad \Longrightarrow \quad
t=\sqrt{\frac{2h}{g}}.
\]

La portée horizontale vaut alors
\[
x_{\text{impact}}=v_0\sqrt{\frac{2h}{g}}.
\]

\subsection{Exemple 4 : virage circulaire}

\begin{exemple}
Une voiture de masse $m$ prend un virage de rayon $R$ à vitesse $v$.
\end{exemple}

Si le mouvement est circulaire uniforme, il faut une accélération centripète :
\[
a=\frac{v^2}{R}.
\]
Donc la résultante des forces latérales doit vérifier
\[
F_{\text{lat}}=m\frac{v^2}{R}.
\]

\begin{interpretation}
Si l'adhérence est insuffisante, la voiture ne peut pas suivre le cercle : elle « part tout droit » tangentiellement.
\end{interpretation}

%====================================================
\section{Clés de compréhension essentielles}
%====================================================

\begin{interpretation}
\textbf{1. Une force ne crée pas forcément de la vitesse, elle crée de l'accélération.}
\end{interpretation}

\begin{interpretation}
\textbf{2. Accélérer ne signifie pas toujours aller plus vite.}  
Tourner à vitesse constante, c'est déjà accélérer.
\end{interpretation}

\begin{interpretation}
\textbf{3. La résultante des forces commande l'évolution du mouvement.}
\end{interpretation}

\begin{interpretation}
\textbf{4. Si la somme des forces est nulle, le mouvement peut être uniforme.}
\end{interpretation}

\begin{interpretation}
\textbf{5. Le poids et la masse sont deux notions différentes.}
\end{interpretation}

\begin{interpretation}
\textbf{6. Une trajectoire courbe nécessite en général une accélération non nulle.}
\end{interpretation}

\begin{interpretation}
\textbf{7. Dans un projectile, l'horizontal et le vertical se traitent séparément.}
\end{interpretation}

\begin{interpretation}
\textbf{8. Les lois de Newton sont des outils de modélisation : elles exigent un système, un référentiel, un bilan des forces et une projection correcte.}
\end{interpretation}

%====================================================
\section{Erreurs classiques à éviter}
%====================================================

\begin{remarque}
Voici les erreurs les plus fréquentes :
\begin{itemize}
    \item confondre vitesse et accélération ;
    \item croire qu'un objet en mouvement subit forcément une force dans le sens du mouvement ;
    \item oublier de préciser le système étudié ;
    \item oublier le référentiel ;
    \item projeter incorrectement une équation vectorielle ;
    \item confondre masse et poids ;
    \item penser que les forces d'action-réaction s'annulent sur le même objet ;
    \item oublier que le mouvement circulaire uniforme est accéléré.
\end{itemize}
\end{remarque}

%====================================================
\section{Ouvertures vers le supérieur}
%====================================================

\subsection{Repère de Frenet}

\begin{remarque}
Dans le supérieur, on décompose souvent l'accélération dans le repère de Frenet :
\[
\vec{a}=a_T\vec{u}_T+a_N\vec{u}_N
\]
avec :
\begin{itemize}
    \item $a_T$ : composante tangentielle ;
    \item $a_N$ : composante normale.
\end{itemize}
\end{remarque}

\subsection{Lois de conservation}

\begin{remarque}
Au-delà de Terminale, on introduit de manière plus systématique :
\begin{itemize}
    \item la conservation de l'énergie mécanique ;
    \item la conservation de la quantité de mouvement ;
    \item le moment cinétique.
\end{itemize}
\end{remarque}

\subsection{Équations différentielles}

\begin{remarque}
Les équations du mouvement
\[
m\vec{a}=\sum \vec{F}
\]
conduisent naturellement à des équations différentielles, très importantes dans le supérieur.
\end{remarque}

%====================================================
\section{Bilan du chapitre}
%====================================================

\begin{center}
\renewcommand{\arraystretch}{1.6}
\begin{tabularx}{\textwidth}{|>{\raggedright\arraybackslash}p{4.6cm}|X|}
\hline
\textbf{Notion} & \textbf{À savoir absolument} \\
\hline
Système & Toujours définir l'objet étudié. \\
\hline
Référentiel & Le mouvement dépend du référentiel choisi. \\
\hline
Vecteur vitesse & \(\vec{v}=\dv{\OM}{t}\), tangent à la trajectoire. \\
\hline
Vecteur accélération & \(\vec{a}=\dv{\vec{v}}{t}\), mesure la variation du vecteur vitesse. \\
\hline
Poids & \(\vec{P}=m\vec{g}\). \\
\hline
Principe d'inertie & Si \(\sum \vec{F}=\vec{0}\), alors repos ou MRU. \\
\hline
Deuxième loi de Newton & \(\sum \vec{F}=m\vec{a}\). \\
\hline
Troisième loi de Newton & \(\vec{F}_{A\to B}=-\vec{F}_{B\to A}\). \\
\hline
Quantité de mouvement & \(\vec{p}=m\vec{v}\). \\
\hline
Chute libre & \(\vec{a}=\vec{g}\). \\
\hline
Projectile & Mouvement horizontal uniforme + vertical uniformément varié. \\
\hline
Mouvement circulaire uniforme & \(a=\frac{v^2}{R}\), dirigée vers le centre. \\
\hline
\end{tabularx}
\end{center}

%====================================================
\section*{Conclusion}
%====================================================

La mécanique de Terminale spécialité permet de comprendre une idée centrale de toute la physique : \textbf{le mouvement ne se résume pas à “aller vite”, il résulte d'interactions modélisées par des forces}.  

Le coeur du chapitre est la relation
\[
\sum \vec{F}=m\vec{a},
\]
qui relie directement les actions mécaniques subies par un système à l'évolution de son mouvement.

C'est un chapitre fondamental, car il apprend à :
\begin{itemize}
    \item modéliser ;
    \item choisir un référentiel ;
    \item faire un bilan des forces ;
    \item écrire des équations ;
    \item interpréter physiquement un résultat.
\end{itemize}

\end{document}